%-----------------------------------------------------------------------------
%  Aerospace-Module-2 - Relazione di progetto (versione italiana)
%
%  Compilazione:  make ita    (oppure: latexmk -pdf main-ita.tex)
%  La versione principale, in inglese, e' main.tex (make).
%  Le figure dei risultati non sono duplicate qui dentro: \graphicspath punta
%  direttamente alle cartelle del progetto che le generano.
%-----------------------------------------------------------------------------
\documentclass[11pt,a4paper]{article}

% --- lingua e font ---------------------------------------------------------
\usepackage[T1]{fontenc}
\usepackage[utf8]{inputenc}
\usepackage{lmodern}
\usepackage[italian]{babel}
\usepackage{microtype}
\usepackage{csquotes}

% --- pagina ----------------------------------------------------------------
\usepackage[margin=2.7cm]{geometry}
\usepackage{fancyhdr}

% --- matematica, tabelle, figure -------------------------------------------
\usepackage{amsmath}
\usepackage{amssymb}
\usepackage{graphicx}
\usepackage{array}
\usepackage{booktabs}
\usepackage{float}
\usepackage{subcaption}
\usepackage{xcolor}
\usepackage{enumitem}

% --- schemi ----------------------------------------------------------------
\usepackage{tikz}
\usetikzlibrary{positioning,arrows.meta}

% --- codice ----------------------------------------------------------------
\usepackage{listings}
\usepackage{upquote}

% --- bibliografia ----------------------------------------------------------
\usepackage[numbers,sort&compress]{natbib}

% --- collegamenti (per ultimo) ---------------------------------------------
\usepackage[hidelinks]{hyperref}

%-----------------------------------------------------------------------------
%  Impostazioni
%-----------------------------------------------------------------------------

% figure prese dove il progetto le produce, senza copiarle nel report
\graphicspath{%
  {immagini/}%
  {../applications/test/nonEqTTv/output/}%
  {../explainations/dipendenze/}%
}

% un po' di elasticita' in piu' nelle righe con nomi lunghi in carattere macchina
\setlength{\emergencystretch}{3em}

\pagestyle{fancy}
\fancyhf{}
\fancyhead[L]{\small\itshape Modello a due temperature in OpenFOAM-13 con Mutation++}
\fancyhead[R]{\small\thepage}
\renewcommand{\headrulewidth}{0.4pt}

\definecolor{codebg}{gray}{0.96}
\definecolor{codekw}{rgb}{0.15,0.30,0.60}
\definecolor{codecm}{rgb}{0.35,0.50,0.35}

\lstdefinestyle{progetto}{%
  backgroundcolor=\color{codebg},
  basicstyle=\ttfamily\footnotesize,
  keywordstyle=\color{codekw}\bfseries,
  commentstyle=\color{codecm}\itshape,
  breaklines=true,
  showstringspaces=false,
  columns=fullflexible,
  keepspaces=true,
  frame=single,
  framerule=0pt,
  xleftmargin=0.6em,
  framexleftmargin=0.6em,
  captionpos=b,
}
\lstset{style=progetto}
\renewcommand{\lstlistingname}{Listato}

% notazione ricorrente
\newcommand{\Ttr}{T_{tr}}
\newcommand{\Tve}{T_{ve}}
\newcommand{\eve}{e_{ve}}
\newcommand{\hyfoam}{hy2Foam}
\newcommand{\mpp}{Mutation\texttt{++}}
\newcommand{\shockThermo}{\texttt{shockThermo}}
\newcommand{\file}[1]{\texttt{#1}}

% segnaposto per le sezioni ancora da scrivere
\newcommand{\daScrivere}{%
  \par\medskip\noindent
  \textcolor{gray}{\itshape [sezione da scrivere]}\par\medskip
}

%-----------------------------------------------------------------------------
\begin{document}

\begin{titlepage}
\centering
\vspace*{2cm}

{\Large Politecnico di Milano\par}
{\large Dipartimento di Scienze e Tecnologie Aerospaziali\par}

\vfill

{\LARGE\bfseries Un modello a due temperature\\[0.2em]
per flussi ipersonici reagenti\\[0.4em]
in OpenFOAM-13 e \mpp\par}

\vfill

{\large Tommaso Marchesini, Federico Quartieri, Daniele Salvi\par}

\vspace*{2cm}
\end{titlepage}

\tableofcontents
\clearpage

\section*{Sommario}
\addcontentsline{toc}{section}{Sommario}

Questa relazione documenta la riproduzione, in OpenFOAM-13 e con la libreria
termochimica \mpp, dei casi di verifica zero-dimensionali con cui Casseau
\textit{et al.}~\cite{casseau2016} hanno validato il solver a due temperature
\hyfoam. Il modello di Park a due temperature -- una temperatura
traslazionale-rotazionale $\Ttr$ e una vibro-elettronica $\Tve$ -- è stato
implementato in forma conservativa, risolvendo le equazioni nelle energie
anziché nelle temperature e delegando a \mpp{} sia l'inversione delle energie
in temperature sia il calcolo delle proprietà termodinamiche.

Il lavoro procede su due percorsi che condividono le stesse funzioni sorgente:
un insieme di programmi zero-dimensionali autonomi, che integrano le equazioni
del bagno termico chiamando soltanto \mpp{} e coprono tutte le figure del
capitolo di verifica del paper, e il solver CFD, che risolve le stesse equazioni su una singola cella
attraverso un modello termodinamico che fa da ponte verso la libreria. L'accordo fra i due percorsi, entro lo 0{,}03--0{,}2\%
su tutte le curve, verifica l'integrazione del modello nel solver; il confronto
con le curve di \hyfoam{} estratte dalle figure vettoriali del paper e con le
temperature di equilibrio ricavate dalla conservazione dell'energia ne verifica
la fisica.


\section{Introduzione}
\label{sec:introduzione}

\subsection{Il non-equilibrio termochimico nel rientro atmosferico}
\label{sec:contesto}

Un veicolo che rientra nell'atmosfera a velocità ipersonica genera un urto
forte, dietro il quale la temperatura sale a decine di migliaia di kelvin: i
modi vibrazionali si eccitano, i livelli elettronici si popolano e le molecole
si dissociano. Questi processi non raggiungono l'equilibrio alla stessa
velocità: traslazione e rotazione si equilibrano in pochi urti molecolari, la
vibrazione in molti di più, e il gas dietro l'urto si trova in non-equilibrio
termico. Il modello a due temperature di Park~\cite{park1990,park1993} lo
descrive tenendo distinte una temperatura traslazionale-rotazionale e una
vibro-elettronica, legate da un tempo di rilassamento. A questo si somma
l'accoppiamento con la chimica: una molecola eccitata vibrazionalmente si
dissocia più facilmente, e le reazioni a loro volta sottraggono o restituiscono
energia vibrazionale. Dalla descrizione di questi processi dipende la
previsione dei carichi termici sul veicolo.

\subsection{Il modello a una temperatura e i suoi limiti}
\label{sec:una-temperatura}

Una molecola può immagazzinare energia in modi diversi: spostandosi nello
spazio (traslazione), ruotando su se stessa (rotazione), facendo oscillare i
propri atomi (vibrazione) ed eccitando i propri elettroni. Nella maggior parte
delle simulazioni CFD, compresi i solver standard di OpenFOAM, il gas è
descritto da una sola temperatura: si assume che l'energia ricevuta dal gas si
ripartisca subito fra tutti questi modi nelle proporzioni di equilibrio, per
cui un solo numero basta a sapere quanta energia c'è in ciascuno, e con lo
stesso numero si calcola anche la velocità delle reazioni chimiche. È
un'ipotesi eccellente nella gran parte delle applicazioni, perché le molecole
si urtano così spesso che la ripartizione avviene molto più in fretta di
quanto cambi il flusso.

Dietro un urto ipersonico questa ipotesi cade: l'energia arriva prima al moto
delle molecole e solo dopo molti urti alla vibrazione, in un tempo paragonabile
a quello che il gas impiega ad attraversare la regione dietro l'urto. Un
modello a una temperatura non vede questo ritardo: prevede appena dietro l'urto
una temperatura più bassa di quella reale e salta la zona in cui il gas si
porta gradualmente all'equilibrio. Lo stesso accade per la chimica: la
dissociazione dipende sia dall'energia con cui le molecole si urtano sia da
quanto stanno già vibrando, due grandezze che una sola temperatura non può
distinguere. Il risultato sono errori sulla temperatura e sulla composizione
del gas, e quindi sulla posizione dell'urto e sul calore che raggiunge la
parete~\cite{park1990}.

\subsection{La soluzione: il modello a due temperature}
\label{sec:soluzione}

Il modello a due temperature di Park~\cite{park1990,park1993} supera questi
limiti usando due temperature invece di una: la prima descrive il moto e la
rotazione delle molecole, che si equilibrano quasi subito, la seconda la
vibrazione e gli elettroni, che restano indietro. Subito dietro l'urto le due
possono differire di migliaia di gradi; l'energia passa poi dall'una all'altra
gradualmente, tanto più in fretta quanto più il gas è caldo e denso, finché le
due temperature coincidono e il gas è tornato in equilibrio. Anche la chimica
le usa entrambe: la dissociazione è calcolata a una temperatura intermedia,
così che una molecola che vibra ancora poco si dissoci più lentamente, e
l'energia di vibrazione portata via dalle molecole che si spezzano viene tolta
da quella del gas. Resta un'approssimazione, perché dentro ciascuno dei due
gruppi l'energia è supposta in equilibrio, ma cattura l'essenziale senza dover
seguire ogni singolo livello energetico delle molecole.

\subsection{Il lavoro di riferimento}
\label{sec:riferimento}

Il lavoro di riferimento è l'articolo di Casseau
\textit{et al.}~\cite{casseau2016}, dell'Università di Strathclyde, che
presenta \hyfoam, un solver a due temperature per flussi ipersonici reagenti
costruito su OpenFOAM e distribuito come codice aperto~\cite{hystrath}. La
prima parte dell'articolo verifica il solver su una serie di bagni termici:
piccoli volumi di gas fermi e isolati, preparati fuori equilibrio e lasciati
evolvere, in cui ogni processo si può studiare da solo. I casi vanno dall'azoto
puro, con e senza energia elettronica, alle miscele di azoto con azoto atomico
e con ossigeno, fino ai gas che reagiscono e all'aria a cinque specie, e i
risultati sono confrontati con altri codici CFD e con codici DSMC, che
simulano il gas seguendo il moto e gli urti di un gran numero di particelle.
Sono questi i casi che il presente lavoro riproduce.


\section{Il lavoro svolto}
\label{sec:lavoro}

\subsection{Approccio generale}
\label{sec:approccio}

Il lavoro mette insieme due programmi che fanno cose diverse.
OpenFOAM~\cite{openfoam13} è la piattaforma CFD: gestisce la mesh, risolve le
equazioni di conservazione cella per cella e fa avanzare la soluzione nel
tempo. \mpp~\cite{scoggins2020} è una libreria di termochimica: dato lo stato
di un gas, calcola quanta energia contiene ciascun modo, quanto dura il
rilassamento vibrazionale e quanto sono veloci le reazioni, e sa fare il
percorso inverso, dalle energie alle temperature. L'idea è lasciare a ciascuno
il proprio mestiere: OpenFOAM risolve le equazioni del modello a due
temperature, \mpp{} fornisce tutte le grandezze termochimiche che vi
compaiono. Il solver, in particolare, non segue direttamente le due
temperature ma le energie corrispondenti, che si conservano correttamente anche
attraverso un urto, e a ogni passo le consegna a \mpp, che ne ricava le
temperature.

Il collegamento con OpenFOAM è realizzato senza modificarne il codice,
attraverso \emph{librerie dinamiche}. Una libreria dinamica è un pezzo di codice
compilato che non fa parte di un programma, ma viene caricato da questo mentre
è in esecuzione, come un'estensione aggiunta a un'applicazione già installata.
OpenFOAM è costruito attorno a questo meccanismo: il solver, il modello
termodinamico, le condizioni al contorno e molti altri componenti non sono
fissati nel programma, ma scelti per nome nei file di testo del caso. Ogni
libreria, quando viene caricata, iscrive i componenti che contiene in un
elenco interno di OpenFOAM; OpenFOAM legge nel caso il nome richiesto, lo cerca
nell'elenco e costruisce il componente corrispondente. Aggiungere un modello
nuovo significa quindi compilare una libreria che lo iscriva nell'elenco
giusto, e OpenFOAM lo userà come uno dei propri.

\subsection{OpenFOAM}
\label{sec:openfoam}

\subsubsection{Com'è fatto e il suo modello a una temperatura}
\label{sec:openfoam-moduli}

OpenFOAM è una piattaforma CFD a codice aperto: un insieme di librerie e di
solver già pronti che discretizzano le equazioni con il metodo dei volumi
finiti, dividendo il dominio in celle e tenendo un valore di ogni grandezza per
cella. È scritta in modo che si possa aggiungere fisica nuova senza toccarne il
nucleo, ed è proprio quello che fa il progetto.

Per i flussi che interessano qui fa già molto. Uno dei suoi solver è fatto
apposta per i flussi compressibili veloci: gestisce le onde d'urto, trasporta
le specie chimiche insieme al flusso e risolve un'equazione per l'energia. A chiudere il
sistema c'è un modello termodinamico: dall'energia di ogni cella ricava una
temperatura, e da quella la pressione, la viscosità, la conducibilità e le
velocità di reazione. Tutta la catena, però, poggia su una sola temperatura: è
il modello a una temperatura della sezione~\ref{sec:una-temperatura}.

\subsubsection{Dove si innesta il progetto}
\label{sec:openfoam-innesto}

Il progetto si innesta in due punti di questa catena. Il suo solver deriva da quel solver per flussi
compressibili veloci, lo stesso da cui è nato \hyfoam, e ne riscrive solo la fase dell'energia: lì aggiunge
l'equazione dell'energia vibro-elettronica e i termini sorgente della chimica,
mentre flussi, densità, quantità di moto e pressione restano quelli ereditati.
Il suo modello termodinamico deriva da quello standard per le miscele
compressibili e aggiunge ciò che serve al modello a due temperature: la
temperatura e l'energia vibro-elettronica come due campi in più, e tre termini
sorgente messi a disposizione del solver. Derivando da quello standard, per il
resto di OpenFOAM è un modello termodinamico come gli altri, e il caso lo
sceglie per nome come qualunque altro. È anche l'unico punto in cui si chiama
\mpp: il solver non parla mai direttamente con la libreria.

\subsection{\mpp}
\label{sec:mutation}

Questa sezione descrive la libreria che fornisce al progetto tutta la
termochimica e il modo in cui il progetto vi si collega.

\subsubsection{Che cos'è e a cosa serve}
\label{sec:mutation-cosa}

\mpp~\cite{scoggins2020} è una libreria scritta in C++, sviluppata
originariamente al von Karman Institute for Fluid Dynamics e distribuita come
codice aperto; il progetto ne usa la versione 1.1.3. È pensata proprio per
essere collegata a un codice CFD: non risolve equazioni di flusso, ma, dato lo
stato di un gas in un punto, ne calcola le proprietà che un solver non sa
calcolare da solo: termodinamiche, di trasporto, chimiche e di scambio di
energia fra i modi, dai flussi subsonici a quelli ipersonici, anche per gas
ionizzati. Il gas è descritto dalle specie che lo compongono, dal numero di
temperature, una oppure due come nel modello di Park, da una banca dati
termodinamica e dall'elenco delle reazioni. Lo stato si può dare in tre modi:
densità e temperature, densità ed energie, oppure frazioni in massa con
pressione e temperature. Da ognuno la libreria ricava tutto il resto, e dalle
energie trova da sé le temperature.

\subsubsection{Il collegamento con il progetto}
\label{sec:mutation-collegamento}

\mpp{} non fa parte di OpenFOAM: è compilata a parte, insieme al progetto, come
libreria dinamica, e il modello termodinamico ad alta entalpia descritto nella
sezione~\ref{sec:openfoam-moduli} e i programmi zero-dimensionali di
verifica ne chiamano le funzioni; il solver ci arriva attraverso il modello. All'avvio di una simulazione il
modello termodinamico legge nel caso la descrizione del gas -- specie, numero
di temperature, banca dati termodinamica, reazioni --, crea il gas
corrispondente in \mpp{} e associa ogni specie di OpenFOAM alla specie di
\mpp{} con lo stesso nome. Da quel momento, cella per cella e a ogni passo, gli
passa lo stato del gas e ne riceve le temperature e tutte le grandezze che
servono al solver. I dati del gas, cioè le costanti delle specie e le
reazioni, sono letti da una cartella del progetto che contiene i valori del
lavoro di riferimento.

\subsection{Il modello implementato}
\label{sec:implementato}

Questa sezione descrive che cosa fa il modello termodinamico ad alta entalpia:
con quali grandezze descrive il gas, quali termini sorgente calcola, come
aggira i limiti di \mpp{} e come il solver lo usa a ogni passo.

\subsubsection{Variabili e stato del gas}
\label{sec:implementato-stato}

In ogni cella il solver non segue le temperature, ma tre grandezze che si
conservano, come nel lavoro di riferimento: la densità di ciascuna specie,
l'energia del gas e la sua parte vibro-elettronica. Le temperature si ricavano
da queste. A ogni passo, dopo che il solver ha aggiornato le energie, il
modello le consegna a \mpp, che usa il modello a due temperature con la banca
dati RRHO e restituisce la temperatura traslazionale e quella
vibro-elettronica; con queste si aggiorna il legame fra pressione e densità.
All'avvio si fa il percorso inverso: dalle temperature assegnate nel caso,
\mpp{} calcola le energie iniziali di ogni cella. Nel passaggio va gestita una
differenza di convenzione: l'energia di OpenFOAM non comprende l'energia
chimica di formazione delle specie, quella di \mpp{} sì, e il modello la
aggiunge o la toglie a seconda della direzione. Il modello a polinomi del
template resta solo per costruire i campi che OpenFOAM si aspetta: tutte le
grandezze del gas vengono da \mpp, e la temperatura iniziale viene riletta dal
caso, perché il modello a polinomi l'avrebbe tagliata a 20\,000~K.

\subsubsection{I termini sorgente}
\label{sec:implementato-sorgenti}

Nel bagno termico le energie e la composizione cambiano solo per effetto dei
termini sorgente, calcolati cella per cella con le funzioni di \mpp. Il primo
è lo scambio fra vibrazione e traslazione, descritto dall'equazione di
Landau--Teller~\cite{landau1936}:
\[
  Q_{VT} = \sum_m \rho_m\,
  \frac{e_{ve,m}(\Ttr) - e_{ve,m}(\Tve)}{\tau_m},
\]
dove la somma è sulle molecole, $\rho_m$ è la densità della molecola $m$,
$e_{ve,m}$ la sua energia vibro-elettronica per unità di massa e $\tau_m$ il
suo tempo di rilassamento, calcolato da \mpp{} con la correlazione di Millikan
e White~\cite{millikan1963} e la correzione di Park. In parole: l'energia
vibrazionale si avvicina al valore che avrebbe alla temperatura
traslazionale, tanto più in fretta quanto più $\tau_m$ è piccolo. Il secondo è
la chimica: \mpp{} fornisce la produzione $\dot\omega_s$ di ogni specie, con le
costanti di Arrhenius del meccanismo valutate alla temperatura di Park
\[
  T_P = \Ttr^{\,a}\,\Tve^{\,1-a}, \qquad a = 0{,}7 .
\]
Le reazioni entrano poi nei bilanci di energia in due modi. Le molecole che si
dissociano portano via la loro energia vibro-elettronica e quelle che si
formano la aggiungono (modello non preferenziale),
\[
  Q_{CV} = \sum_s e_{ve,s}(\Tve)\,\dot\omega_s ;
\]
inoltre, poiché l'energia di OpenFOAM non comprende quella di formazione,
l'energia spesa per rompere i legami chimici va tolta a quella del gas con il
calore di reazione $Q_{ch} = -\sum_s h_{f,s}\,\dot\omega_s$, dove $h_{f,s}$ è
l'entalpia di formazione della specie $s$. In sintesi: le specie cambiano per
$\dot\omega_s$, l'energia vibro-elettronica per $Q_{VT} + Q_{CV}$ e l'energia
del gas, che la comprende, per $Q_{ch}$, così che l'energia totale, compresa
quella chimica, si conserva.

\subsubsection{Come sono stati aggirati i limiti di \mpp}
\label{sec:implementato-limiti}

I modelli di \mpp{} non coincidono in tutto con quelli del lavoro di
riferimento. Le differenze, e il modo in cui sono state trattate, sono queste:

\begin{itemize}[leftmargin=1.6em]
  \item \textbf{Energia che guida lo scambio V--T.} \mpp{} usa la sola
    energia vibrazionale, il lavoro di riferimento anche quella elettronica
    della molecola. Lo scambio è stato riscritto come nel lavoro di
    riferimento, con le energie e i tempi di rilassamento della libreria;
    senza livelli elettronici le due forme coincidono.
  \item \textbf{Temperatura delle dissociazioni.} \mpp{} usa sempre
    $\sqrt{\Ttr\Tve}$, cioè $a = 0{,}5$, e non si può cambiare. Per ottenere
    $a = 0{,}7$ si porta per un istante la libreria in uno stato in cui
    entrambe le temperature valgono $T_P$, si chiedono le velocità di reazione
    e si ripristina lo stato vero. Il procedimento vale solo per meccanismi
    fatti di sole dissociazioni, come quello dell'azoto usato qui.
  \item \textbf{Tempo di rilassamento nelle miscele.} \mpp{} media i tempi
    relativi ai vari partner di collisione in modo diverso dal lavoro di
    riferimento. Si è tenuta la versione della libreria, per non riscriverne
    il modello; l'effetto è discusso insieme ai risultati.
  \item \textbf{Una sola temperatura vibrazionale.} Nei casi risolti con il
    solver l'unica molecola presente è l'azoto, e il limite non conta. Il caso
    con azoto e ossigeno è risolto solo con un programma zero-dimensionale,
    che tiene separate le due vibrazioni e calcola a parte lo scambio V--V fra
    di esse.
  \item \textbf{Accoppiamento chimica-vibrazione.} \mpp{} offre solo il
    modello non preferenziale, ed è quello usato; il modello CVDV del lavoro
    di riferimento non è stato riprodotto.
  \item \textbf{Reazioni inverse.} Il set di reazioni per l'aria contiene due
    coppie di reazioni irreversibili una inversa dell'altra, che \mpp{} non
    accetta nello stesso meccanismo: sono state divise in due meccanismi,
    le cui velocità di reazione vengono sommate.
\end{itemize}

\subsubsection{Il solver}
\label{sec:implementato-solver}

Il solver a due temperature è derivato dal modulo density-based di OpenFOAM:
ne eredita il calcolo dei flussi, della densità, della quantità di moto e della
pressione, e riscrive la fase dell'energia. In questa fase, a ogni passo:

\begin{enumerate}[leftmargin=1.8em]
  \item risolve le equazioni delle specie, con la produzione chimica fornita
    da \mpp, e corregge le frazioni perché la loro somma resti uno;
  \item risolve l'equazione dell'energia del gas, uguale a quella del modulo
    originale più il calore di reazione;
  \item risolve l'equazione dell'energia vibro-elettronica, che per ora
    contiene solo la variazione nel tempo e i termini sorgente: manca il
    trasporto con il flusso, che nel bagno termico non serve;
  \item consegna le energie a \mpp, che restituisce le due temperature, e
    aggiorna il legame fra pressione e densità.
\end{enumerate}

\noindent I termini sorgente sono calcolati con lo stato del gas già noto e
trattati in modo esplicito. La fase dell'energia del modulo originale non
poteva essere riusata perché, subito dopo aver risolto l'equazione
dell'energia, ne ricava la temperatura: mancherebbero sia il calore di
reazione sia l'equazione vibro-elettronica, che deve essere risolta prima di
convertire le energie in temperature.

\subsubsection{Le equazioni risolte: forma generale e caso del bagno termico}
\label{sec:implementato-equazioni}

La Sezione~\ref{sec:implementato-sorgenti} ha dato i termini sorgente; questa
sezione riporta le equazioni di conservazione in cui entrano, prima nella
forma generale per cui \texttt{shockFluid} e \shockThermo{} sono scritti, valida
su una mesh qualunque, poi nella forma ridotta che viene effettivamente
integrata nei casi a bagno termico usati per la verifica
(Sezione~\ref{sec:programmi0D} e i casi \shockThermo{} corrispondenti).

Su una mesh generica, massa e quantità di moto sono risolte senza modifiche
rispetto al modulo originale \texttt{shockFluid},
\[
  \frac{\partial \rho}{\partial t} + \nabla\cdot(\rho\mathbf{U}) = 0 ,
  \qquad
  \frac{\partial (\rho\mathbf{U})}{\partial t}
  + \nabla\cdot(\rho\mathbf{U}\mathbf{U}) + \nabla p = \nabla\cdot\tau ,
\]
con lo sforzo viscoso newtoniano $\tau = \mu\left[\nabla\mathbf{U} +
(\nabla\mathbf{U})^{T} - \tfrac{2}{3}(\nabla\cdot\mathbf{U})\mathbf{I}\right]$,
dove $\mu$ è la viscosità dinamica. L'equazione dell'energia risolta da
\shockThermo{} aggiunge il calore di reazione a quella del modulo originale,
\[
  \frac{\partial (\rho e)}{\partial t} + \nabla\cdot(\boldsymbol\phi_p)
  + \frac{\partial (\rho K)}{\partial t}
  = Q_{ch} + \nabla\cdot(\tau\cdot\mathbf{U}) + \nabla\cdot(\kappa\nabla T) ,
\]
dove $K = \tfrac{1}{2}|\mathbf{U}|^2$ è l'energia cinetica per unità di massa,
$\kappa$ la conducibilità termica, e $\boldsymbol\phi_p$ il flusso convettivo
di energia attraverso ciascuna faccia, comprensivo del lavoro di pressione,
ricostruito con lo stesso schema di Kurganov--Tadmor usato per massa e
quantità di moto. Per analogia, una forma generale (non a bagno termico)
dell'equazione vibro-elettronica aggiungerebbe un termine convettivo e uno
diffusivo dello stesso tipo,
\[
  \frac{\partial (\rho\eve)}{\partial t} + \nabla\cdot(\phi\,\eve)
  = Q_{VT} + Q_{CV} + \nabla\cdot(\kappa_{ve}\nabla\Tve) ,
\]
ma, come notato nella Sezione~\ref{sec:implementato-solver}, sono
attualmente implementati solo la variazione nel tempo e i termini sorgente.

Nessuno dei termini di trasporto qui sopra è mai attraversato dai casi
risolti in questo progetto. Ogni caso è un bagno termico: una singola cella
di calcolo, con campo uniforme nello spazio, ogni condizione al contorno a
gradiente nullo e il gas fermo, $\mathbf{U}=0$.
In queste condizioni $\nabla\mathbf{U}=0$ e $\nabla T=0$ identicamente, in
ogni cella e a ogni passo temporale, quindi
\[
  \tau = 0, \qquad \kappa\nabla T = 0, \qquad K = 0, \qquad
  \boldsymbol\phi_p = 0 ,
\]
indipendentemente dal valore numerico di $\mu$ e $\kappa$; massa e quantità
di moto si riducono all'identità banale $0=0$. Quello che resta, ed è ciò
che viene effettivamente integrato, è
\[
  \frac{\partial (\rho e)}{\partial t} = Q_{ch} ,
  \qquad
  \frac{\partial (\rho \eve)}{\partial t} = Q_{VT} + Q_{CV} .
\]
Sono le stesse due equazioni differenziali ordinarie avanzate nel tempo dai
programmi 0D della Sezione~\ref{sec:programmi0D}, ed è questo che rende il
confronto fra le due implementazioni un vero controllo del solver, e non
della mesh o della discretizzazione.

\subsection{I programmi di verifica zero-dimensionali}
\label{sec:programmi0D}

Accanto al solver il progetto comprende quattro piccoli programmi che simulano
il bagno termico senza OpenFOAM, uno per ogni famiglia di casi del lavoro di
riferimento: l'azoto puro, la miscela di azoto molecolare e atomico, con o
senza reazioni, la miscela di azoto e ossigeno e l'aria a cinque specie.
Ognuno prepara il gas nello stato iniziale della figura da riprodurre, calcola
le energie contenute nel volume e avanza nel tempo con passi di un
nanosecondo, come il lavoro di riferimento: a ogni passo calcola i termini
sorgente, aggiorna l'energia vibro-elettronica e la composizione, e chiede a
\mpp{} le nuove temperature. L'energia totale, compresa quella chimica, resta
costante, perché il volume è chiuso e isolato. I termini sorgente sono
calcolati con le stesse funzioni usate dal solver, ed è questo che rende
significativo il confronto fra i due percorsi. Nei casi senza reazioni i
programmi calcolano anche la temperatura finale attesa, cercando la
temperatura unica alla quale il gas avrebbe la stessa energia dello stato
iniziale: è il valore verso cui le due temperature devono convergere.

Due programmi si discostano da questo schema. Quello per azoto e ossigeno tiene
separate le energie vibrazionali delle due molecole, ciascuna con la propria
temperatura, e aggiunge lo scambio V--V fra di esse con la formulazione di
Knab \textit{et al.}~\cite{knab1995} usata dal lavoro di riferimento e la
probabilità di scambio consigliata di $10^{-2}$; la sezione d'urto della coppia
azoto-ossigeno, che il lavoro di riferimento non riporta, è presa da
\hyfoam~\cite{hystrath}. Quello per l'aria usa invece una sola temperatura,
come fa il lavoro di riferimento per questo caso. I programmi sono anche molto
più rapidi del solver: il caso più lungo, un milione di passi, richiede circa
un secondo contro quasi mezz'ora.

\subsection{I dati del gas}
\label{sec:dati}

Le proprietà delle specie usate da \mpp{} sono quelle delle tabelle del lavoro
di riferimento: temperature caratteristiche di vibrazione, entalpie di
formazione e livelli elettronici delle cinque specie dell'aria. I parametri del
rilassamento vibrazionale sono quelli già contenuti in \mpp, tratti da
Park~\cite{park1993}, che per azoto e ossigeno coincidono con quelli del lavoro
di riferimento.

L'energia elettronica richiede una scelta. Il lavoro di riferimento la tratta
esplicitamente solo nel caso dell'azoto a 30\,000~K, dove confronta il gas con
e senza; per gli altri casi non lo dice, e il testo che commenta la miscela di
azoto molecolare e atomico parla anzi del contributo elettronico dell'azoto
atomico. Le temperature finali riportate nell'articolo, però, tornano solo
senza livelli elettronici: con i livelli elettronici, la temperatura finale
della miscela di azoto molecolare e atomico sarebbe di circa 18\,300~K invece
dei 24\,400~K dell'articolo, e quella della miscela di azoto e ossigeno di
circa 19\,200~K invece di 12\,100~K, mentre senza si ottengono 24\,353~K e
12\,133~K. Si sono quindi seguiti i risultati: il progetto ha due versioni dei
dati, una con i livelli elettronici, usata solo per il caso dell'azoto che li
include, e una con il solo livello fondamentale, usata in tutti gli altri
casi.

Le reazioni sono di due tipi. Per la dissociazione dell'azoto si usa la
reazione irreversibile $\mathrm{N_2 + N_2 \rightarrow 2\,N + N_2}$ con le
costanti di Park riportate nel lavoro di riferimento. Per l'aria a cinque
specie il lavoro di riferimento usa diciannove reazioni irreversibili, quindici
dissociazioni e quattro scambi, con costanti ricavate dalla teoria
Quantum-Kinetic~\cite{scanlon2015}, che cita ma non riporta: sono state prese
dai dati di \hyfoam~\cite{hystrath}. Per confronto è stato usato anche il
meccanismo per l'aria già presente in \mpp, con le costanti di Park.

% ---------------------------------------------------------------------------
% Traccia per le prossime sottosezioni (dalla vecchia sezione "Il modello"):
%  - decomposizione dell'energia nei modi (eq. 1-7 del paper) e le due T
%  - variabili conservate: rho_s, E, E_ve (eq. 23); perche' energie e non T
%  - equazione di stato: p = somma delle pressioni parziali (eq. 24)
%  - scambio V-T: Landau-Teller (eq. 8), tau di Millikan-White + Park (eq. 9-17)
%  - scambio V-V: Knab (eq. 18), solo nel caso N2-O2
%  - chimica: legge di azione di massa (eq. 27), Arrhenius (eq. 28),
%    temperatura di Park T^a Tv^(1-a) (eq. 29)
%  - accoppiamento chimica-vibrazione: modello non preferenziale (eq. 30-31)
%  - [fatto, vedi sec:implementato-equazioni] riepilogo del sistema risolto
%    nel bagno termico (eq. 22, 25, 26)


\section{Verifica e risultati}
\label{sec:risultati}

\subsection{Metodo}
\label{sec:metodo}

Ogni caso del lavoro di riferimento è stato calcolato con il programma
zero-dimensionale corrispondente; i casi delle fig.~3a, 3b, 4, 5 e 7 sono
stati calcolati anche con il solver, su una mesh di una sola cella, con le
stesse condizioni iniziali, gli stessi dati e lo stesso passo di un
nanosecondo. I risultati sono sottoposti a tre verifiche:

\begin{enumerate}[leftmargin=1.8em]
  \item \textbf{solver contro programmi zero-dimensionali}: i due percorsi
    fanno la stessa fisica, quindi ogni differenza viene dal modo in cui il
    solver integra il modello; è la verifica dell'implementazione;
  \item \textbf{programmi zero-dimensionali contro il lavoro di riferimento}:
    è la verifica della fisica;
  \item \textbf{temperatura finale contro conservazione dell'energia}: la
    temperatura raggiunta all'equilibrio deve essere quella calcolata
    dall'energia iniziale, e quella indicata nel lavoro di riferimento.
\end{enumerate}

Le curve di \hyfoam{} sono state ricavate direttamente dal pdf dell'articolo.
Le figure sono disegni vettoriali: ogni curva è memorizzata come una sequenza
di punti, che è stata letta e riportata in tempo e temperatura usando le tacche
degli assi. L'incertezza è circa lo spessore della linea, 40--60~K sulle
temperature. Gli scarti sulle temperature sono dati in percentuale, come
massimo lungo il transitorio. Le densità delle specie sono disegnate in scala
logaritmica, e il loro scarto è dato in \emph{decadi}: una decade è un fattore
10, e 0{,}01 decadi corrispondono a circa il 2{,}3\%. I programmi salvano i
risultati ogni 10 o 100~ns, e i confronti partono da lì. Nei grafici le linee
chiare continue sono il lavoro di riferimento (\emph{paper}), quelle
tratteggiate i programmi zero-dimensionali (\emph{standalone}) e i punti il
solver (\emph{solver}).

\subsection{Azoto puro}
\label{sec:risultati-azoto}

Le fig.~3a e 3b del lavoro di riferimento riguardano l'azoto a 1~atm, riscaldato
nella vibrazione (10\,000~K per la traslazione, 1000~K per la vibrazione) o
raffreddato (3000~K e 10\,000~K). Nel riscaldamento la temperatura finale è
7623{,}3~K, esattamente il valore dato nel testo dell'articolo, e le curve
differiscono al massimo dello 0{,}3\% sulla temperatura traslazionale e
dell'1{,}5\% su quella vibrazionale. Nel raffreddamento la temperatura finale è
4973~K contro 4972~K; lo scarto nel transitorio sale al 2{,}8\% e al 5{,}1\%,
e corrisponde a un piccolo ritardo del rilassamento
(figura~\ref{fig:azoto-3}).

\begin{figure}[ht]
  \centering
  \begin{subfigure}{0.49\textwidth}
    \includegraphics[width=\linewidth]{fig3a}
    \caption{Riscaldamento vibrazionale (fig.~3a).}
  \end{subfigure}\hfill
  \begin{subfigure}{0.49\textwidth}
    \includegraphics[width=\linewidth]{fig3b}
    \caption{Raffreddamento vibrazionale (fig.~3b).}
  \end{subfigure}
  \caption{Azoto puro a 1~atm: temperatura traslazionale ($T_{tr}$) e
  vibrazionale ($T_v$) nel tempo.}
  \label{fig:azoto-3}
\end{figure}

La fig.~4 parte da 30\,000~K e mette a confronto il caso senza e con energia
elettronica (figura~\ref{fig:azoto-4}). Le temperature finali sono 21\,932~K
e 17\,650~K, contro i 21{,}9 e 17{,}7~kK indicati in figura: l'energia
elettronica abbassa la temperatura finale di oltre 4000~K. Gli scarti nel
transitorio sono dello 0{,}05\% e dello 0{,}8\% senza energia elettronica,
dello 0{,}3\% e del 4\% con. Questo caso giustifica la scelta di riscrivere lo
scambio V--T con l'energia vibro-elettronica
(sezione~\ref{sec:implementato-limiti}): con la forma originale di \mpp, che
usa la sola energia vibrazionale, la temperatura finale sarebbe la stessa ma
gli scarti nel transitorio salirebbero al 28\% e al 56\%.

\begin{figure}[ht]
  \centering
  \begin{subfigure}{0.49\textwidth}
    \includegraphics[width=\linewidth]{fig4-noEl}
    \caption{Senza energia elettronica.}
  \end{subfigure}\hfill
  \begin{subfigure}{0.49\textwidth}
    \includegraphics[width=\linewidth]{fig4-el}
    \caption{Con energia elettronica.}
  \end{subfigure}
  \caption{Azoto puro da 30\,000~K (fig.~4); la linea punteggiata orizzontale
  è la temperatura finale attesa dalla conservazione dell'energia.}
  \label{fig:azoto-4}
\end{figure}

\subsection{Miscele non reagenti}
\label{sec:risultati-miscele}

La fig.~5 riguarda una miscela di azoto molecolare e atomico, con
$5\cdot10^{22}$ particelle per metro cubo di ciascuna specie, a 30\,000~K e
1000~K. La temperatura finale è 24\,353~K contro i 24{,}4~kK della figura, e la
temperatura traslazionale resta entro il 2{,}5\%; quella vibrazionale, però,
sale più tardi che nel lavoro di riferimento, con uno scarto che arriva al
20\% nel tratto più ripido. L'origine più probabile è il tempo di
rilassamento nelle miscele, che \mpp{} media in modo diverso dal lavoro di
riferimento (sezione~\ref{sec:implementato-limiti}): il rilassamento
vibrazionale risulta un po' più lento, mentre lo stato finale è corretto.

La fig.~6 riguarda una miscela di azoto e ossigeno in parti uguali, a 1~atm,
con la vibrazione molto più calda della traslazione (30\,000~K contro 5000~K),
calcolata senza e con lo scambio V--V fra le due molecole. La temperatura
finale è 12\,133~K contro i 12{,}1~kK della figura. Senza scambio V--V gli
scarti sono del 2{,}7\% sulla temperatura traslazionale, del 2{,}7\% su quella
vibrazionale dell'azoto e del 6{,}4\% su quella dell'ossigeno; con lo scambio
scendono al 2{,}0\%, 2{,}5\% e 3{,}4\%. Come nel lavoro di riferimento, lo
scambio V--V avvicina le due temperature vibrazionali e rende più rapido il
ritorno all'equilibrio (figura~\ref{fig:miscele}).

\begin{figure}[ht]
  \centering
  \begin{subfigure}{0.49\textwidth}
    \includegraphics[width=\linewidth]{fig5}
    \caption{Azoto molecolare e atomico (fig.~5).}
  \end{subfigure}\hfill
  \begin{subfigure}{0.49\textwidth}
    \includegraphics[width=\linewidth]{fig6-VV}
    \caption{Azoto e ossigeno, con scambio V--V (fig.~6).}
  \end{subfigure}
  \caption{Miscele non reagenti. Nella miscela di azoto e ossigeno ogni
  molecola ha la propria temperatura vibrazionale.}
  \label{fig:miscele}
\end{figure}

\subsection{Miscele reagenti}
\label{sec:risultati-reagenti}

Le fig.~7 e 8 riprendono la miscela di azoto molecolare e atomico, ora con la
dissociazione dell'azoto. Nella fig.~7, che parte da 30\,000~K e 1000~K, alla
fine del calcolo le due temperature valgono 8309~K e 8312~K contro 8292~K e
8281~K del lavoro di riferimento, e le densità dell'azoto molecolare e atomico,
riferite a quella iniziale, valgono 0{,}138 e 1{,}225 contro 0{,}139 e 1{,}222:
la composizione segue quella del lavoro di riferimento entro 0{,}006 decadi,
cioè l'1{,}4\%, per tutto il calcolo (figura~\ref{fig:reagenti}). La
temperatura traslazionale resta entro il 2{,}3\%; quella vibrazionale mostra
lo stesso ritardo della fig.~5, fino al 17\% all'inizio. La fig.~8 parte
invece da 30\,000~K per entrambe le temperature; le temperature finali sono
12\,121~K e 12\,141~K contro 12\,101~K e 12\,043~K, la composizione è entro
0{,}01 decadi e le temperature entro il 2\% e il 10\%.

Questi casi confermano anche la scelta dell'esponente della temperatura di
Park. Con il valore fisso di \mpp{} ($a = 0{,}5$) la composizione della fig.~7
si allontana dal lavoro di riferimento fino a 0{,}038 decadi (9\%) e la
temperatura traslazionale fino al 6{,}7\%; con $a = 0{,}7$, come nel lavoro di
riferimento, gli scarti scendono a 0{,}006 decadi e al 2{,}3\%. La temperatura
vibrazionale non cambia, perché il suo scarto viene dal rilassamento e non
dalla chimica.

\begin{figure}[ht]
  \centering
  \begin{subfigure}{0.49\textwidth}
    \includegraphics[width=\linewidth]{fig7}
    \caption{Temperature.}
  \end{subfigure}\hfill
  \begin{subfigure}{0.49\textwidth}
    \includegraphics[width=\linewidth]{fig7-n}
    \caption{Densità delle specie riferite a quella iniziale.}
  \end{subfigure}
  \caption{Miscela reagente di azoto molecolare e atomico da 30\,000~K e
  1000~K (fig.~7), con esponente di Park $a = 0{,}7$.}
  \label{fig:reagenti}
\end{figure}

\subsection{Aria a cinque specie}
\label{sec:risultati-aria}

La fig.~9 riguarda aria (79\% di azoto e 21\% di ossigeno in moli) a
0{,}063~atm e 10\,000~K, con diciannove reazioni e una sola temperatura. Con le
costanti Quantum-Kinetic la temperatura segue il lavoro di riferimento entro lo
0{,}7\% (5096~K contro 5074~K alla fine) e le densità delle cinque specie entro
0{,}06 decadi, quattro su cinque entro 0{,}03 (figura~\ref{fig:aria}). Con il
meccanismo di Park di \mpp{}, invece, la temperatura si scosta fino all'8{,}9\%
e le densità di ossigeno molecolare e atomico, azoto atomico e monossido di
azoto di 0{,}2--0{,}5 decadi, cioè di un fattore fra 1{,}6 e 3: la fig.~9 del
lavoro di riferimento è stata ottenuta con le costanti Quantum-Kinetic.

\begin{figure}[ht]
  \centering
  \begin{subfigure}{0.49\textwidth}
    \includegraphics[width=\linewidth]{fig9-QK}
    \caption{Temperatura.}
  \end{subfigure}\hfill
  \begin{subfigure}{0.49\textwidth}
    \includegraphics[width=\linewidth]{fig9-QK-n}
    \caption{Densità delle specie riferite a quella iniziale.}
  \end{subfigure}
  \caption{Aria a cinque specie da 10\,000~K con le costanti Quantum-Kinetic
  (fig.~9).}
  \label{fig:aria}
\end{figure}

\subsection{Solver e programmi zero-dimensionali}
\label{sec:risultati-solver}

La tabella~\ref{tab:solver} riporta lo scarto massimo fra il solver e i
programmi zero-dimensionali in tutti i casi calcolati con entrambi. Le
temperature finali coincidono, e le curve restano entro lo 0{,}2\% in tutti i
casi tranne la fig.~4, dove lo scarto arriva all'1{,}2\% nel rilassamento
molto rapido dei primi 100~ns. Lì conta di più il modo in cui si avanza nel
tempo: il solver ripete il calcolo due volte a ogni passo, i programmi
zero-dimensionali una sola. Anche con le reazioni accese la composizione
coincide entro lo 0{,}03\%. Il modello è quindi integrato correttamente nel
solver, chimica compresa.

\begin{table}[ht]
  \centering
  \caption{Scarto massimo fra solver e programmi zero-dimensionali.}
  \label{tab:solver}
  \small
  \begin{tabular}{lccc}
    \toprule
    Caso & $T_{tr}$ & $T_{v}$ & Composizione \\
    \midrule
    fig.~3b, raffreddamento            & 0{,}00\% & 0{,}00\% & -- \\
    fig.~3a, riscaldamento             & 0{,}03\% & 0{,}21\% & -- \\
    fig.~4, senza energia elettronica  & 0{,}14\% & 0{,}69\% & -- \\
    fig.~4, con energia elettronica    & 0{,}51\% & 1{,}18\% & -- \\
    fig.~5, azoto molecolare e atomico & 0{,}02\% & 0{,}13\% & -- \\
    fig.~7, con reazioni               & 0{,}03\% & 0{,}18\% & 0{,}03\% \\
    \bottomrule
  \end{tabular}
\end{table}


\section{Discussione}
\label{sec:discussione}

\subsection{Che cosa dicono i risultati}
\label{sec:discussione-risultati}

Il confronto fra solver e programmi zero-dimensionali mostra che il modello è
stato inserito correttamente nel solver: le due strade, che fanno la stessa
fisica, danno le stesse curve entro lo 0{,}2\% in quasi tutti i casi, anche
con le reazioni accese. Il confronto con il lavoro di riferimento mostra che
anche la fisica è quella giusta. Nei casi senza reazioni le temperature finali
coincidono con quelle imposte dalla conservazione dell'energia e con i valori
del lavoro di riferimento entro lo 0{,}3\%, cioè entro l'incertezza con cui si
leggono le figure; nei casi con reazioni le temperature a fine calcolo
coincidono entro lo 0{,}8\%, la composizione della miscela di azoto entro il
2{,}5\% e quella dell'aria entro il 15\%. Nel transitorio l'accordo
è di pochi punti percentuali per l'azoto puro, per la miscela di azoto e
ossigeno e per l'aria. La differenza più grande riguarda la temperatura
vibrazionale nelle miscele di azoto molecolare e atomico, dove arriva al
10--20\%.

\subsection{L'origine delle differenze}
\label{sec:discussione-origine}

L'origine più probabile del ritardo della temperatura vibrazionale nelle
miscele è il modo in cui \mpp{} media i tempi di rilassamento relativi ai
diversi partner di collisione, che non è quello del lavoro di riferimento. Il
modo di mediare conta solo quando i partner sono più di uno: nell'azoto puro
l'accordo è buono, mentre nelle miscele con azoto atomico la temperatura
vibrazionale sale più tardi. Lo stato finale non ne risente, perché dipende
solo dall'energia. L'ipotesi non è stata verificata implementando la media del
lavoro di riferimento, che resta uno degli sviluppi possibili. Per l'azoto puro
e per la miscela di azoto e ossigeno restano piccoli spostamenti nel tempo, di
qualche punto percentuale, la cui origine non è stata isolata.

\subsection{Che cosa insegna il confronto}
\label{sec:discussione-lezioni}

Buona parte del lavoro è consistita nel ricostruire scelte che l'articolo non
dichiara, o dichiara in modo incompleto: se l'energia elettronica è accesa,
quale energia guida lo scambio vibrazionale, con quale esponente si calcola la
temperatura delle dissociazioni, quali costanti di velocità sono usate per
l'aria. Nessuna di queste scelte è un dettaglio: sbagliarne una sposta i
risultati molto più delle differenze rimaste, come mostrano i confronti della
sezione~\ref{sec:risultati}, con scarti fino al 56\% sulla temperatura
vibrazionale, un fattore 3 sulle densità dell'aria e circa il 25\% sulla
temperatura finale. Lo stesso articolo ne dà un esempio quando spiega la
differenza con LeMANS, nel caso della miscela di azoto molecolare e atomico,
con una diversa convenzione nella
correzione di Park.

Per prendere queste decisioni due strumenti sono stati decisivi. Il primo è
la temperatura finale imposta dalla conservazione dell'energia, che non
dipende da nessun modello di rilassamento e dice subito se i dati del gas sono
quelli giusti. Il secondo sono le curve estratte dalle figure vettoriali, che
trasformano il confronto con l'articolo da una lettura a occhio in una misura.
Infine, la verifica su due strade separa gli errori di implementazione dalle
scelte di modello: poiché solver e programmi coincidono, ogni differenza
rimasta rispetto al lavoro di riferimento viene dal modello, non da un errore
nel codice.


\section{Conclusioni e sviluppi futuri}
\label{sec:conclusioni}

\subsection{Conclusioni}
\label{sec:conclusioni-sintesi}

Il lavoro ha portato in OpenFOAM-13 un modello a due temperature per flussi
ipersonici reagenti, appoggiato alla libreria \mpp{} per tutta la
termochimica. Il solver risolve le equazioni nelle energie, che si conservano
anche attraverso un urto, e ricava le temperature con \mpp; le funzioni
sorgente sono scritte una sola volta e condivise con programmi
zero-dimensionali indipendenti. Dove i modelli della libreria non coincidevano
con quelli del lavoro di riferimento, le sorgenti sono state riscritte sopra di
essa.

Tutti i casi di verifica zero-dimensionali di Casseau
\textit{et al.}~\cite{casseau2016} sono stati riprodotti. Il solver coincide
con i programmi zero-dimensionali entro lo 0{,}3\% in tutti i casi tranne
l'azoto ad alta temperatura, dove arriva all'1{,}2\%, reazioni comprese: il modello è integrato correttamente.
Rispetto al lavoro di riferimento gli stati finali coincidono entro l'1\%, la
composizione della miscela di azoto entro il 2{,}5\% e quella dell'aria entro
il 15\%; l'unica differenza rilevante è un ritardo della temperatura
vibrazionale nelle miscele con azoto atomico, con ogni probabilità dovuto al
modo in cui \mpp{} media i tempi di rilassamento. Il confronto ha anche
mostrato quanto i risultati dipendano da scelte che il lavoro di riferimento non dichiara,
come l'uso dell'energia elettronica o le costanti di velocità per l'aria, e ha
permesso di ricostruirle. Il modello è verificato sul bagno termico, dove non
c'è flusso: per usarlo in casi reali servono gli sviluppi che seguono.

\subsection{Sviluppi futuri}
\label{sec:sviluppi}

Gli sviluppi proposti riguardano il calcolo su più celle, il trasporto
dell'energia e i flussi viscosi. Il primo è quello decisivo e si appoggia agli
altri due.

\begin{enumerate}[leftmargin=1.8em]
  \item \textbf{Calcolo su una griglia vera.} Finora tutto è stato verificato
    su una sola cella: è questo il passo che rende il modello utilizzabile su
    un flusso vero. Che cosa serve è descritto nella
    sezione~\ref{sec:discussione-multicella}: condizioni al contorno per i due
    campi nuovi, sorgenti calcolate con la densità del passo corrente e i
    termini di trasporto del punto successivo. Per il resto non ci sono
    ostacoli, perché tutti i calcoli del modello termodinamico sono locali
    alla cella.
  \item \textbf{Trasporto dell'energia vibro-elettronica.} È la parte del
    punto precedente che si può sviluppare e provare per conto suo. I flussi
    alle facce che servono per il trasporto sono quelli già calcolati per le
    specie, mentre la conduzione usa la conducibilità vibro-elettronica del
    punto successivo.
  \item \textbf{Viscosità e conducibilità a due temperature.} Servono per i
    flussi viscosi e per il calore che raggiunge la parete. \mpp{} le calcola
    già, e divide la conducibilità fra la parte traslazionale e quella
    vibro-elettronica: basta chiederle cella per cella nel modello
    termodinamico, come oggi le energie, e passarle al solver al posto di
    quelle del modello a una temperatura.
\end{enumerate}


\clearpage

\bibliographystyle{unsrtnat}
\bibliography{bibliografia}

\end{document}
